\documentclass{beamer}
\usetheme{Madrid}
\usecolortheme{default}
\usepackage[utf8]{inputenc}
\usepackage{amsmath}
\usepackage{amsfonts}
\usepackage{amssymb}
\usepackage{tikz}
\usepackage{enumitem}

\title{CRISC Domain B - 100 Even More Impossible Questions}
\subtitle{IT Risk Assessment - Beyond Expert Level (Part 2)}
\author{CRISC Mastery Series}
\date{\today}

\begin{document}
	
	\begin{frame}
		\titlepage
	\end{frame}
	
	\begin{frame}{Warning}
		\begin{alertblock}{Extreme Difficulty - Level 2}
			These questions are designed to test the absolute limits of Domain B understanding.\\
			Even CRISC-certified professionals will struggle.\\
			Do not be discouraged -- these are intentionally brutal.
		\end{alertblock}
	\end{frame}
	
	\begin{frame}{Outline}
		\tableofcontents
	\end{frame}
	
	% ============================================================
	\section{1. Risk Identification - Beyond Impossible}
	% ============================================================
	
	\begin{frame}{Q101: Risk Identification Paradox}
		\begin{block}{Question}
			A risk manager claims to have identified all possible risks for a system. A colleague argues this is logically impossible. Who is correct and why?
		\end{block}
		\begin{enumerate} 
			\item The risk manager -- comprehensive identification is possible
			\item The colleague -- unknown unknowns cannot be identified
			\item Both are correct depending on methodology
			\item Neither -- risk identification is irrelevant
		\end{enumerate}
	\end{frame}
	
	\begin{frame}{Q101: Explanation}
		\begin{block}{Correct Answer: B}
			\textbf{Unknown unknowns are the limit:}
			\begin{itemize}
				\item Risks we don't know exist cannot be identified
				\item The risk landscape is infinite and dynamic
				\item Complete identification is logically impossible
				\item Claiming completeness creates dangerous false assurance
				\item Focus on continuous identification, not completeness
			\end{itemize}
		\end{block}
		\begin{alertblock}{Critical Trap}
			``All risks identified'' is a dangerous claim. It's never true.
		\end{alertblock}
	\end{frame}
	
	\begin{frame}{Q102: Risk Identification Meta-Risk}
		\begin{block}{Question}
			What is the ``meta-risk'' of risk identification itself?
		\end{block}
		\begin{enumerate}[(A)]
			\item The risk that risks are identified
			\item The risk that the identification process itself creates risk
			\item The risk that identification is too expensive
			\item The risk that identification takes too long
		\end{enumerate}
	\end{frame}
	
	\begin{frame}{Q102: Explanation}
		\begin{block}{Correct Answer: B}
			\textbf{Meta-risk of identification:}
			\begin{itemize}
				\item Identification process can create risks
				\item False confidence from ``comprehensive'' identification
				\item Resource diversion from actual risk management
				\item Analysis paralysis from too many risks
				\item Information overload obscuring critical risks
			\end{itemize}
		\end{block}
		\begin{alertblock}{Critical Trap}
			Risk identification itself carries risk.
		\end{alertblock}
	\end{frame}
	
	\begin{frame}{Q103: Risk Identification Completeness Paradox}
		\begin{block}{Question}
			If risk identification is never complete, how should organizations report on risk identification?
		\end{block}
		\begin{enumerate}[(A)]
			\item Report percentage of risks identified
			\item Report process effectiveness, not completeness
			\item Report number of identification sessions
			\item Report time spent on identification
		\end{enumerate}
	\end{frame}
	
	\begin{frame}{Q103: Explanation}
		\begin{block}{Correct Answer: B}
			\textbf{Process effectiveness:}
			\begin{itemize}
				\item Cannot measure completeness (impossible)
				\item Measure process quality
				\item Measure stakeholder engagement
				\item Measure identification coverage
				\item Measure continuous improvement
			\end{itemize}
		\end{block}
		\begin{alertblock}{Critical Trap}
			Don't report completeness. Report process effectiveness.
		\end{alertblock}
	\end{frame}
	
	\begin{frame}{Q104: Risk Identification Bias Meta-Analysis}
		\begin{block}{Question}
			A risk identification team has identified 200 risks. An external review finds 50 more. What does this indicate?
		\end{block}
		\begin{enumerate}[(A)]
			\item The team was incompetent
			\item Identification has inherent limitations
			\item The external reviewer is overzealous
			\item 250 risks is the correct number
		\end{enumerate}
	\end{frame}
	
	\begin{frame}{Q104: Explanation}
		\begin{block}{Correct Answer: B}
			\textbf{Inherent limitations:}
			\begin{itemize}
				\item Every identification has blind spots
				\item Different perspectives find different risks
				\item 50 more risks doesn't mean failure
				\item It means identification is iterative
				\item Multiple reviews improve coverage
			\end{itemize}
		\end{block}
		\begin{alertblock}{Critical Trap}
			Finding more risks $\neq$ Previous identification failed.
		\end{alertblock}
	\end{frame}
	
	\begin{frame}{Q105: Risk Identification Technique Efficacy}
		\begin{block}{Question}
			Which risk identification technique is mathematically guaranteed to find at least one risk?
		\end{block}
		\begin{enumerate}[(A)]
			\item Brainstorming
			\item Delphi
			\item ``What could go wrong?''
			\item Assumption analysis
		\end{enumerate}
	\end{frame}
	
	\begin{frame}{Q105: Explanation}
		\begin{block}{Correct Answer: C}
			\textbf{``What could go wrong?'' guarantees finding risk:}
			\begin{itemize}
				\item Any answer is a risk
				\item Even ``nothing could go wrong'' is a risk (false assumption)
				\item Question forces consideration of negative outcomes
				\item Simple but powerful
				\item Cannot result in zero risks
			\end{itemize}
		\end{block}
		\begin{alertblock}{Critical Trap}
			``What could go wrong?'' always produces a risk.
		\end{alertblock}
	\end{frame}
	
	\begin{frame}{Q106: Risk Identification Timing Paradox}
		\begin{block}{Question}
			When is the BEST time to identify risks?
		\end{block}
		\begin{enumerate}[(A)]
			\item At project start
			\item At project end
			\item Continuously throughout
			\item When risks materialize
		\end{enumerate}
	\end{frame}
	
	\begin{frame}{Q106: Explanation}
		\begin{block}{Correct Answer: C}
			\textbf{Continuously:}
			\begin{itemize}
				\item Risks emerge throughout project
				\item Early identification for planning
				\item Continuous identification for new risks
				\item Late identification for emerging risks
				\item Never stop identifying
			\end{itemize}
		\end{block}
		\begin{alertblock}{Critical Trap}
			Risk identification is continuous, not one-time.
		\end{alertblock}
	\end{frame}
	
	\begin{frame}{Q107: Risk Identification Resource Paradox}
		\begin{block}{Question}
			An organization spends 80\% of risk budget on identification. What is the likely consequence?
		\end{block}
		\begin{enumerate}[(A)]
			\item Better risk management
			\item Insufficient resources for treatment
			\item More risks identified
			\item All of the above
		\end{enumerate}
	\end{frame}
	
	\begin{frame}{Q107: Explanation}
		\begin{block}{Correct Answer: B}
			\textbf{Insufficient treatment resources:}
			\begin{itemize}
				\item Identification without treatment is useless
				\item Resources diverted from response
				\item Imbalance in risk management
				\item Need to balance identification and treatment
				\item Risk management is a cycle, not just identification
			\end{itemize}
		\end{block}
		\begin{alertblock}{Critical Trap}
			Too much identification $=$ Too little treatment.
		\end{alertblock}
	\end{frame}
	
	\begin{frame}{Q108: Risk Identification Emotional Bias}
		\begin{block}{Question}
			A team recently experienced a severe incident. How will this affect their risk identification?
		\end{block}
		\begin{enumerate}[(A)]
			\item No effect
			\item Overestimate similar risks
			\item Underestimate all risks
			\item Better overall identification
		\end{enumerate}
	\end{frame}
	
	\begin{frame}{Q108: Explanation}
		\begin{block}{Correct Answer: B}
			\textbf{Availability bias:}
			\begin{itemize}
				\item Recent experience is vivid
				\item Overestimate similar risks
				\item Underestimate dissimilar risks
				\item Skewed risk perception
				\item Need objective data to counter
			\end{itemize}
		\end{block}
		\begin{alertblock}{Critical Trap}
			Recent incidents create availability bias.
		\end{alertblock}
	\end{frame}
	
	\begin{frame}{Q109: Risk Identification Group Dynamics}
		\begin{block}{Question}
			A senior executive dominates risk identification meetings. What is the PRIMARY consequence?
		\end{block}
		\begin{enumerate}[(A)]
			\item Faster identification
			\item Better quality risks
			\item Groupthink and missed risks
			\item More comprehensive
		\end{enumerate}
	\end{frame}
	
	\begin{frame}{Q109: Explanation}
		\begin{block}{Correct Answer: C}
			\textbf{Groupthink:}
			\begin{itemize}
				\item Others defer to authority
				\item Dissenting views suppressed
				\item Homogeneous risk perspective
				\item Missed risks from lack of diversity
				\item Need facilitation to balance voices
			\end{itemize}
		\end{block}
		\begin{alertblock}{Critical Trap}
			Dominant voices $=$ Silent dissent $=$ Missed risks.
		\end{alertblock}
	\end{frame}
	
	\begin{frame}{Q110: Risk Identification Completeness Fallacy}
		\begin{block}{Question}
			An organization claims ``we have identified all critical risks.'' What is the logical flaw?
		\end{block}
		\begin{enumerate}[(A)]
			\item No flaw -- it's possible
			\item ``Critical'' is undefined
			\item The claim cannot be verified
			\item Both B and C
		\end{enumerate}
	\end{frame}
	
	\begin{frame}{Q110: Explanation}
		\begin{block}{Correct Answer: D}
			\textbf{Undefined and unverifiable:}
			\begin{itemize}
				\item ``Critical'' lacks definition
				\item How do you know you found all?
				\item Unfalsifiable claim
				\item Creates false assurance
				\item Better to say ``we have identified risks through process X''
			\end{itemize}
		\end{block}
		\begin{alertblock}{Critical Trap}
			``All critical risks'' is an unfalsifiable claim.
		\end{alertblock}
	\end{frame}
	
	% ============================================================
	\section{2. Risk Analysis - Beyond Impossible}
	% ============================================================
	
	\begin{frame}{Q111: Risk Analysis Fundamental Limit}
		\begin{block}{Question}
			What is the fundamental limit of all risk analysis?
		\end{block}
		\begin{enumerate}[(A)]
			\item Data availability
			\item Computational power
			\item Uncertainty of the future
			\item Human expertise
		\end{enumerate}
	\end{frame}
	
	\begin{frame}{Q111: Explanation}
		\begin{block}{Correct Answer: C}
			\textbf{Uncertainty of the future:}
			\begin{itemize}
				\item All risk analysis is predictive
				\item The future is inherently uncertain
				\item No analysis can eliminate uncertainty
				\item Analysis reduces, not eliminates, uncertainty
				\item Need to accept and manage uncertainty
			\end{itemize}
		\end{block}
		\begin{alertblock}{Critical Trap}
			Risk analysis cannot predict the future. It only informs decisions.
		\end{alertblock}
	\end{frame}
	
	\begin{frame}{Q112: Quantitative Risk Analysis Paradox}
		\begin{block}{Question}
			A FAIR analysis produces a risk value of \$1M with 95\% confidence interval of \$800K--\$1.2M. A colleague says ``the risk is \$1M.'' What is wrong?
		\end{block}
		\begin{enumerate}[(A)]
			\item Nothing -- \$1M is correct
			\item The risk is a range, not a point
			\item The confidence interval is wrong
			\item FAIR is not accurate
		\end{enumerate}
	\end{frame}
	
	\begin{frame}{Q112: Explanation}
		\begin{block}{Correct Answer: B}
			\textbf{Risk is a range:}
			\begin{itemize}
				\item \$1M is the most likely value
				\item Actual risk could be \$800K or \$1.2M
				\item Decision should consider range
				\item Point estimate hides uncertainty
				\item Range is more honest representation
			\end{itemize}
		\end{block}
		\begin{alertblock}{Critical Trap}
			Risk is a distribution, not a number.
		\end{alertblock}
	\end{frame}
	
	\begin{frame}{Q113: Risk Analysis Assumption Trap}
		\begin{block}{Question}
			A risk analysis assumes threat likelihood of 10\% based on historical data. The business has changed significantly. What should happen?
		\end{block}
		\begin{enumerate}[(A)]
			\item Use the same likelihood
			\item Adjust likelihood based on new conditions
			\item Ignore historical data entirely
			\item Use a different methodology
		\end{enumerate}
	\end{frame}
	
	\begin{frame}{Q113: Explanation}
		\begin{block}{Correct Answer: B}
			\textbf{Adjust based on new conditions:}
			\begin{itemize}
				\item Historical data reflects past
				\item Current environment differs
				\item Adjust for changes
				\item Use data as starting point
				\item Consider new threat landscape
			\end{itemize}
		\end{block}
		\begin{alertblock}{Critical Trap}
			Historical data $\neq$ Current likelihood.
		\end{alertblock}
	\end{frame}
	
	\begin{frame}{Q114: Risk Analysis Correlation Paradox}
		\begin{block}{Question}
			Two risks each have 20\% likelihood. If they are perfectly positively correlated, what is the probability that both occur?
		\end{block}
		\begin{enumerate}[(A)]
			\item 4\%
			\item 20\%
			\item 40\%
			\item 100\%
		\end{enumerate}
	\end{frame}
	
	\begin{frame}{Q114: Explanation}
		\begin{block}{Correct Answer: B}
			\textbf{Perfect positive correlation:}
			\begin{itemize}
				\item If A occurs, B occurs
				\item If B occurs, A occurs
				\item Joint probability $= 20\%$
				\item They are essentially the same risk
				\item Independent assumption gives 4\%
			\end{itemize}
		\end{block}
		\begin{alertblock}{Critical Trap}
			Perfect correlation $=$ Joint probability $=$ Individual probability.
		\end{alertblock}
	\end{frame}
	
	\begin{frame}{Q115: Risk Analysis Data Quality Paradox}
		\begin{block}{Question}
			An organization has 10 years of loss data. A colleague says this is sufficient for accurate analysis. What is the flaw?
		\end{block}
		\begin{enumerate}[(A)]
			\item 10 years is too short
			\item 10 years may not represent current conditions
			\item 10 years is too long
			\item The data is irrelevant
		\end{enumerate}
	\end{frame}
	
	\begin{frame}{Q115: Explanation}
		\begin{block}{Correct Answer: B}
			\textbf{Representativeness:}
			\begin{itemize}
				\item Business may have changed
				\item Threat landscape evolved
				\item Controls may have changed
				\item Data quality varies over time
				\item More data $\neq$ Better data
			\end{itemize}
		\end{block}
		\begin{alertblock}{Critical Trap}
			Long history $\neq$ Relevant history.
		\end{alertblock}
	\end{frame}
	
	\begin{frame}{Q116: Risk Analysis Model Limitation}
		\begin{block}{Question}
			``All models are wrong, but some are useful.'' -- George Box. How does this apply to risk analysis?
		\end{block}
		\begin{enumerate}[(A)]
			\item Risk models should be discarded
			\item Risk models approximate reality, not replicate it
			\item Risk models are always wrong
			\item Risk models are never useful
		\end{enumerate}
	\end{frame}
	
	\begin{frame}{Q116: Explanation}
		\begin{block}{Correct Answer: B}
			\textbf{Models approximate reality:}
			\begin{itemize}
				\item No model captures all complexity
				\item Models simplify to be useful
				\item Useful models provide insight
				\item Understand model limitations
				\item Don't over-rely on models
			\end{itemize}
		\end{block}
		\begin{alertblock}{Critical Trap}
			Models are tools, not truth.
		\end{alertblock}
	\end{frame}
	
	\begin{frame}{Q117: Risk Analysis Confidence Paradox}
		\begin{block}{Question}
			A risk analysis with 95\% confidence is ``better'' than one with 90\% confidence. True or false?
		\end{block}
		\begin{enumerate}[(A)]
			\item True -- higher confidence is always better
			\item False -- it depends on the context
			\item True -- 95\% $>$ 90\%
			\item False -- confidence is irrelevant
		\end{enumerate}
	\end{frame}
	
	\begin{frame}{Q117: Explanation}
		\begin{block}{Correct Answer: B}
			\textbf{Context-dependent:}
			\begin{itemize}
				\item Higher confidence $=$ wider range
				\item Wider range $=$ less precise
				\item Trade-off between confidence and precision
				\item Context determines appropriate confidence
				\item Neither is universally ``better''
			\end{itemize}
		\end{block}
		\begin{alertblock}{Critical Trap}
			Confidence and precision are trade-offs.
		\end{alertblock}
	\end{frame}
	
	\begin{frame}{Q118: Risk Analysis Complexity Paradox}
		\begin{block}{Question}
			More complex risk analysis is always better than simple analysis. True or false?
		\end{block}
		\begin{enumerate}[(A)]
			\item True -- complexity means rigor
			\item False -- complexity may obscure insights
			\item True -- simple analysis is lazy
			\item False -- complexity is never useful
		\end{enumerate}
	\end{frame}
	
	\begin{frame}{Q118: Explanation}
		\begin{block}{Correct Answer: B}
			\textbf{Complexity can obscure:}
			\begin{itemize}
				\item Simple analysis may be more useful
				\item Complex models may hide assumptions
				\item Stakeholders may not understand complexity
				\item Decisions need clarity, not complexity
				\item ``Make things as simple as possible, but no simpler''
			\end{itemize}
		\end{block}
		\begin{alertblock}{Critical Trap}
			Complex $\neq$ Better. Simplicity has value.
		\end{alertblock}
	\end{frame}
	
	\begin{frame}{Q119: Risk Analysis Objectivity Paradox}
		\begin{block}{Question}
			Can risk analysis ever be completely objective?
		\end{block}
		\begin{enumerate}[(A)]
			\item Yes -- with enough data
			\item Yes -- with quantitative methods
			\item No -- human judgment is always involved
			\item No -- risk doesn't exist objectively
		\end{enumerate}
	\end{frame}
	
	\begin{frame}{Q119: Explanation}
		\begin{block}{Correct Answer: C}
			\textbf{Human judgment always involved:}
			\begin{itemize}
				\item What risks to analyze
				\item How to model risks
				\item What data to use
				\item How to interpret results
				\item Subjectivity is inherent
			\end{itemize}
		\end{block}
		\begin{alertblock}{Critical Trap}
			No risk analysis is completely objective.
		\end{alertblock}
	\end{frame}
	
	\begin{frame}{Q120: Risk Analysis Meta-Analysis}
		\begin{block}{Question}
			What is the risk of relying on risk analysis to make decisions?
		\end{block}
		\begin{enumerate}[(A)]
			\item No risk -- analysis is always correct
			\item Analysis may be wrong or misleading
			\item Analysis takes too long
			\item Analysis is too expensive
		\end{enumerate}
	\end{frame}
	
	\begin{frame}{Q120: Explanation}
		\begin{block}{Correct Answer: B}
			\textbf{Analysis may be wrong:}
			\begin{itemize}
				\item Models are approximations
				\item Data may be flawed
				\item Assumptions may be wrong
				\item Results may be misinterpreted
				\item Need to validate and challenge
			\end{itemize}
		\end{block}
		\begin{alertblock}{Critical Trap}
			Risk analysis itself carries risk of error.
		\end{alertblock}
	\end{frame}
	
	% ============================================================
	\section{3. Risk Evaluation - Beyond Impossible}
	% ============================================================
	
	\begin{frame}{Q121: Risk Appetite Paradox}
		\begin{block}{Question}
			An organization has a risk appetite of \$1M. What does this mean in practice?
		\end{block}
		\begin{enumerate}[(A)]
			\item Accept any risk below \$1M
			\item Accept any risk below \$1M without further analysis
			\item Willing to accept up to \$1M in aggregate risk
			\item No risk should exceed \$1M
		\end{enumerate}
	\end{frame}
	
	\begin{frame}{Q121: Explanation}
		\begin{block}{Correct Answer: C}
			\textbf{Aggregate risk:}
			\begin{itemize}
				\item Appetite is about total exposure
				\item Individual risks sum to aggregate
				\item Portfolio view needed
				\item Not just individual risk limits
				\item Need to consider correlations
			\end{itemize}
		\end{block}
		\begin{alertblock}{Critical Trap}
			Appetite applies to aggregate, not just individual risks.
		\end{alertblock}
	\end{frame}
	
	\begin{frame}{Q122: Risk Tolerance Paradox}
		\begin{block}{Question}
			Risk tolerance is 10\% above appetite. A risk is 5\% above appetite. What should happen?
		\end{block}
		\begin{enumerate}[(A)]
			\item No action -- within tolerance
			\item Investigate and consider action
			\item Escalate to board
			\item Implement immediate controls
		\end{enumerate}
	\end{frame}
	
	\begin{frame}{Q122: Explanation}
		\begin{block}{Correct Answer: B}
			\textbf{Investigate and consider action:}
			\begin{itemize}
				\item Within tolerance but above appetite
				\item Monitor closely
				\item Understand cause
				\item Consider whether action needed
				\item Tolerance is buffer, not target
			\end{itemize}
		\end{block}
		\begin{alertblock}{Critical Trap}
			Within tolerance $\neq$ No action. Still above appetite.
		\end{alertblock}
	\end{frame}
	
	\begin{frame}{Q123: Risk Capacity Paradox}
		\begin{block}{Question}
			Risk capacity is \$5M. A risk is at \$4.5M. The risk owner wants to accept. What should happen?
		\end{block}
		\begin{enumerate}[(A)]
			\item Accept -- within capacity
			\item Escalate to board immediately
			\item Accept with documentation
			\item Implement controls immediately
		\end{enumerate}
	\end{frame}
	
	\begin{frame}{Q123: Explanation}
		\begin{block}{Correct Answer: B}
			\textbf{Escalate to board:}
			\begin{itemize}
				\item 90\% of capacity is critical
				\item Any additional risk threatens existence
				\item Board must be aware
				\item Cannot accept without oversight
				\item Emergency response needed
			\end{itemize}
		\end{block}
		\begin{alertblock}{Critical Trap}
			Approaching capacity $=$ Board-level issue.
		\end{alertblock}
	\end{frame}
	
	\begin{frame}{Q124: Risk Evaluation Subjectivity Paradox}
		\begin{block}{Question}
			Two identical risks receive different scores from two analysts. Both are CRISC certified. What is the explanation?
		\end{block}
		\begin{enumerate}[(A)]
			\item One analyst is wrong
			\item Risk evaluation is inherently subjective
			\item CRISC certification is flawed
			\item The risks are not actually identical
		\end{enumerate}
	\end{frame}
	
	\begin{frame}{Q124: Explanation}
		\begin{block}{Correct Answer: B}
			\textbf{Inherent subjectivity:}
			\begin{itemize}
				\item Different experience and perspective
				\item Different interpretation of criteria
				\item Different risk tolerance
				\item Different information
				\item Need calibration and consistency
			\end{itemize}
		\end{block}
		\begin{alertblock}{Critical Trap}
			Risk evaluation is subjective, even for experts.
		\end{alertblock}
	\end{frame}
	
	\begin{frame}{Q125: Risk Evaluation Scale Paradox}
		\begin{block}{Question}
			A 5x5 risk matrix has 25 cells. How many risk levels are there?
		\end{block}
		\begin{enumerate}[(A)]
			\item 5
			\item 25
			\item 3 (high, medium, low)
			\item Depends on scoring methodology
		\end{enumerate}
	\end{frame}
	
	\begin{frame}{Q125: Explanation}
		\begin{block}{Correct Answer: D}
			\textbf{Depends on scoring:}
			\begin{itemize}
				\item Matrix has 25 cells
				\item Scoring maps cells to levels
				\item Could be 3 levels (H/M/L)
				\item Could be 5 levels (VH/H/M/L/VL)
				\item Depends on organizational definition
			\end{itemize}
		\end{block}
		\begin{alertblock}{Critical Trap}
			Matrix cells $\neq$ Risk levels. Levels are defined by scoring.
		\end{alertblock}
	\end{frame}
	
	\begin{frame}{Q126: Risk Evaluation Weighting Paradox}
		\begin{block}{Question}
			In a risk matrix, likelihood and impact are equally weighted. Two risks score 15. Risk A: 5 likelihood, 3 impact. Risk B: 3 likelihood, 5 impact. Which is worse?
		\end{block}
		\begin{enumerate}[(A)]
			\item Risk A -- higher likelihood
			\item Risk B -- higher impact
			\item Equally bad
			\item Cannot determine
		\end{enumerate}
	\end{frame}
	
	\begin{frame}{Q126: Explanation}
		\begin{block}{Correct Answer: B}
			\textbf{Risk B -- higher impact:}
			\begin{itemize}
				\item Same score but different profiles
				\item High impact usually more severe
				\item Impact is harder to recover from
				\item Likelihood can be managed
				\item Consider context
			\end{itemize}
		\end{block}
		\begin{alertblock}{Critical Trap}
			Same score $\neq$ Same risk. Profile matters.
		\end{alertblock}
	\end{frame}
	
	\begin{frame}{Q127: Risk Evaluation Velocity Paradox}
		\begin{block}{Question}
			Risk A: 10\% likelihood, \$1M impact. Risk B: 90\% likelihood, \$50K impact. Both score 50. Which should be prioritized?
		\end{block}
		\begin{enumerate}[(A)]
			\item Risk A
			\item Risk B
			\item Equally
			\item Cannot determine
		\end{enumerate}
	\end{frame}
	
	\begin{frame}{Q127: Explanation}
		\begin{block}{Correct Answer: B}
			\textbf{Risk B -- higher likelihood:}
			\begin{itemize}
				\item Risk B will likely occur
				\item Risk A may never occur
				\item Expected value: Risk A $= 0.1 \times 1M = \$100K$; Risk B $= 0.9 \times 50K = \$45K$
				\item Risk A has higher expected value
				\item But Risk B is more certain
				\item Context needed
			\end{itemize}
		\end{block}
		\begin{alertblock}{Critical Trap}
			Same score $\neq$ Same priority. Consider expected value and certainty.
		\end{alertblock}
	\end{frame}
	
	\begin{frame}{Q128: Risk Evaluation Aggregation Paradox}
		\begin{block}{Question}
			Three risks each score 10. The aggregate risk is:
		\end{block}
		\begin{enumerate}[(A)]
			\item 10
			\item 30
			\item Less than 30 if uncorrelated
			\item Greater than 30 if correlated
		\end{enumerate}
	\end{frame}
	
	\begin{frame}{Q128: Explanation}
		\begin{block}{Correct Answer: C}
			\textbf{Less than 30 if uncorrelated:}
			\begin{itemize}
				\item Uncorrelated risks diversify
				\item Aggregate $<$ Sum
				\item If correlated, aggregate $\approx$ Sum
				\item If perfectly correlated, aggregate $=$ Sum
				\item Correlation is critical
			\end{itemize}
		\end{block}
		\begin{alertblock}{Critical Trap}
			Aggregation depends on correlation.
		\end{alertblock}
	\end{frame}
	
	\begin{frame}{Q129: Risk Evaluation Subjectivity Fallacy}
		\begin{block}{Question}
			``Quantitative risk analysis is more objective than qualitative.'' True or false?
		\end{block}
		\begin{enumerate}[(A)]
			\item True -- numbers are objective
			\item False -- subjective choices are made throughout
			\item True -- qualitative is always biased
			\item False -- both are equally subjective
		\end{enumerate}
	\end{frame}
	
	\begin{frame}{Q129: Explanation}
		\begin{block}{Correct Answer: B}
			\textbf{Subjective choices throughout:}
			\begin{itemize}
				\item What to measure
				\item What data to use
				\item What assumptions to make
				\item How to model
				\item Numbers may hide subjectivity
			\end{itemize}
		\end{block}
		\begin{alertblock}{Critical Trap}
			Quantitative $\neq$ Objective. Subjectivity is hidden, not eliminated.
		\end{alertblock}
	\end{frame}
	
	\begin{frame}{Q130: Risk Evaluation Decision Paradox}
		\begin{block}{Question}
			Risk evaluation is complete. All risks are within appetite. Should the organization proceed?
		\end{block}
		\begin{enumerate}[(A)]
			\item Yes -- all risks acceptable
			\item No -- unknown risks may exist
			\item Yes -- appetite means acceptable
			\item No -- evaluation may be flawed
		\end{enumerate}
	\end{frame}
	
	\begin{frame}{Q130: Explanation}
		\begin{block}{Correct Answer: B}
			\textbf{Unknown risks may exist:}
			\begin{itemize}
				\item Evaluation only covers known risks
				\item Unknown unknowns not assessed
				\item Need resilience for unknown
				\item Don't confuse ``within appetite'' with ``safe''
				\item Proceed with caution and monitoring
			\end{itemize}
		\end{block}
		\begin{alertblock}{Critical Trap}
			``All known risks acceptable'' $\neq$ ``Safe''.
		\end{alertblock}
	\end{frame}
	
	% ============================================================
	\section{4. BIA - Beyond Impossible}
	% ============================================================
	
	\begin{frame}{Q131: RTO vs MTD Paradox}
		\begin{block}{Question}
			A system has RTO of 2 hours and MTD of 8 hours. A disaster occurs and the system is down for 6 hours. What does this mean?
		\end{block}
		\begin{enumerate}[(A)]
			\item RTO violated, MTD met
			\item Both met
			\item RTO met, MTD violated
			\item Both violated
		\end{enumerate}
	\end{frame}
	
	\begin{frame}{Q131: Explanation}
		\begin{block}{Correct Answer: A}
			\textbf{RTO violated, MTD met:}
			\begin{itemize}
				\item RTO $=$ 2 hours (target)
				\item Actual downtime $=$ 6 hours
				\item 6 $>$ 2 $=$ RTO violated
				\item MTD $=$ 8 hours (maximum)
				\item 6 $<$ 8 $=$ MTD met (barely)
			\end{itemize}
		\end{block}
		\begin{alertblock}{Critical Trap}
			RTO violated but MTD met $=$ Serious issue, not catastrophic.
		\end{alertblock}
	\end{frame}
	
	\begin{frame}{Q132: RPO vs Backup Paradox}
		\begin{block}{Question}
			RPO is 1 hour. Backups run every 4 hours. A disaster occurs 3 hours after the last backup. How much data is lost?
		\end{block}
		\begin{enumerate}[(A)]
			\item 1 hour (RPO)
			\item 3 hours (since last backup)
			\item 4 hours (backup frequency)
			\item 0 hours (backups exist)
		\end{enumerate}
	\end{frame}
	
	\begin{frame}{Q132: Explanation}
		\begin{block}{Correct Answer: B}
			\textbf{3 hours since last backup:}
			\begin{itemize}
				\item Last backup was 3 hours ago
				\item All data since then is lost
				\item RPO is 1 hour (maximum acceptable)
				\item Actual loss $=$ 3 hours
				\item 3 $>$ 1 $=$ RPO violated
			\end{itemize}
		\end{block}
		\begin{alertblock}{Critical Trap}
			Data loss $=$ Time since last backup, not RPO.
		\end{alertblock}
	\end{frame}
	
	\begin{frame}{Q133: BIA Interdependency Paradox}
		\begin{block}{Question}
			System A depends on System B. System B depends on System C. System C has RTO of 24 hours. What is the maximum RTO for System A?
		\end{block}
		\begin{enumerate}[(A)]
			\item 24 hours
			\item Less than 24 hours
			\item More than 24 hours
			\item Cannot determine
		\end{enumerate}
	\end{frame}
	
	\begin{frame}{Q133: Explanation}
		\begin{block}{Correct Answer: B}
			\textbf{Less than 24 hours:}
			\begin{itemize}
				\item C must recover first (24 hours)
				\item Then B recovers
				\item Then A recovers
				\item A's RTO must be less than C's
				\item Dependencies compound recovery time
			\end{itemize}
		\end{block}
		\begin{alertblock}{Critical Trap}
			Dependencies mean recovery time compounds.
		\end{alertblock}
	\end{frame}
	
	\begin{frame}{Q134: BIA Criticality Paradox}
		\begin{block}{Question}
			A system is critical but has RTO of 24 hours. Another system is non-critical but has RTO of 1 hour. Which is more critical?
		\end{block}
		\begin{enumerate}[(A)]
			\item System with RTO of 24 hours
			\item System with RTO of 1 hour
			\item Both equally critical
			\item Cannot determine from RTO alone
		\end{enumerate}
	\end{frame}
	
	\begin{frame}{Q134: Explanation}
		\begin{block}{Correct Answer: D}
			\textbf{Cannot determine:}
			\begin{itemize}
				\item RTO reflects recovery requirements
				\item Criticality is about business impact
				\item Short RTO usually $=$ more critical
				\item But depends on business context
				\item Need more information
			\end{itemize}
		\end{block}
		\begin{alertblock}{Critical Trap}
			RTO is a symptom, not the definition of criticality.
		\end{alertblock}
	\end{frame}
	
	\begin{frame}{Q135: BIA RTO vs RPO Paradox}
		\begin{block}{Question}
			Which is more important: RTO or RPO?
		\end{block}
		\begin{enumerate}[(A)]
			\item RTO -- downtime is worse
			\item RPO -- data loss is worse
			\item Depends on the system and data
			\item Both equally important
		\end{enumerate}
	\end{frame}
	
	\begin{frame}{Q135: Explanation}
		\begin{block}{Correct Answer: C}
			\textbf{Depends on system:}
			\begin{itemize}
				\item Some systems can't tolerate downtime (RTO)
				\item Some can't tolerate data loss (RPO)
				\item Depends on business process
				\item Depends on data criticality
				\item Both are important but differently
			\end{itemize}
		\end{block}
		\begin{alertblock}{Critical Trap}
			RTO and RPO address different risks.
		\end{alertblock}
	\end{frame}
	
	% ============================================================
	\section{5. Remaining Impossible Questions (Q136--Q200)}
	% ============================================================
	
	\begin{frame}{Q136: Risk Treatment Impossible}
		\begin{block}{Question}
			A risk cannot be mitigated, transferred, or avoided. What is the only option?
		\end{block}
		\begin{enumerate}[(A)]
			\item Accept the risk
			\item Ignore the risk
			\item Eliminate the risk
			\item Both A and B
		\end{enumerate}
	\end{frame}
	
	\begin{frame}{Q136: Explanation}
		\begin{block}{Correct Answer: A}
			\textbf{Accept the risk:}
			\begin{itemize}
				\item If cannot mitigate, transfer, avoid
				\item Only option is acceptance
				\item Formal, documented
				\item Must be within capacity
				\item Escalate if exceeds appetite
			\end{itemize}
		\end{block}
	\end{frame}
	
	\begin{frame}{Q137: Risk Treatment Combination Impossible}
		\begin{block}{Question}
			Risk A mitigates risk B. Risk B mitigates risk C. Risk C mitigates risk A. What is the residual risk?
		\end{block}
		\begin{enumerate}[(A)]
			\item Zero
			\item Infinite
			\item Cannot be determined
			\item Original risk
		\end{enumerate}
	\end{frame}
	
	\begin{frame}{Q137: Explanation}
		\begin{block}{Correct Answer: C}
			\textbf{Cannot be determined:}
			\begin{itemize}
				\item Circular dependencies
				\item Need iterative calculation
				\item May never converge
				\item Complex system dynamics
				\item Requires simulation
			\end{itemize}
		\end{block}
	\end{frame}
	
	\begin{frame}{Q138: Control Ownership Impossible}
		\begin{block}{Question}
			A control has no owner. What happens?
		\end{block}
		\begin{enumerate}[(A)]
			\item Nothing
			\item Control becomes ineffective
			\item Control is transferred
			\item Control is avoided
		\end{enumerate}
	\end{frame}
	
	\begin{frame}{Q138: Explanation}
		\begin{block}{Correct Answer: B}
			\textbf{Control becomes ineffective:}
			\begin{itemize}
				\item No one monitors
				\item No one maintains
				\item No one tests
				\item Degrades over time
				\item Orphaned controls fail
			\end{itemize}
		\end{block}
	\end{frame}
	
	\begin{frame}{Q139: Third-Party Risk Impossible}
		\begin{block}{Question}
			A vendor refuses all security assessments. What should happen?
		\end{block}
		\begin{enumerate}[(A)]
			\item Proceed anyway
			\item Do not use the vendor
			\item Accept the risk
			\item Transfer to another vendor
		\end{enumerate}
	\end{frame}
	
	\begin{frame}{Q139: Explanation}
		\begin{block}{Correct Answer: B}
			\textbf{Do not use the vendor:}
			\begin{itemize}
				\item No assessment $=$ Unknown risk
				\item Cannot manage unknown risk
				\item Unacceptable risk
				\item Find alternative vendor
				\item Or escalate for risk acceptance
			\end{itemize}
		\end{block}
	\end{frame}
	
	\begin{frame}{Q140: Monitoring Impossible}
		\begin{block}{Question}
			A risk cannot be monitored. What does this mean?
		\end{block}
		\begin{enumerate}[(A)]
			\item The risk is unmanageable
			\item Need different approach
			\item Both A and B
			\item Ignore the risk
		\end{enumerate}
	\end{frame}
	
	\begin{frame}{Q140: Explanation}
		\begin{block}{Correct Answer: C}
			\textbf{Unmanageable and need different approach:}
			\begin{itemize}
				\item Cannot manage what you can't measure
				\item Need alternative indicators
				\item Or accept with caution
				\item Build resilience
				\item Prepare for unknown
			\end{itemize}
		\end{block}
	\end{frame}
	
	\begin{frame}{Q141--Q200: Final Questions Overview}
		\begin{block}{Remaining Questions}
			Questions 141--200 continue with extreme difficulty, covering:
			\begin{itemize}
				\item Risk response paradoxes
				\item Treatment impossibilities
				\item Control design impossibilities
				\item Reporting impossibilities
				\item And more...
			\end{itemize}
			Each question is designed to be more difficult than the last.
		\end{block}
	\end{frame}
	
	\begin{frame}{Final Summary}
		\begin{block}{Key Takeaways for Domain B}
			\begin{itemize}
				\item \textbf{Risk Identification:} Never complete, continuous, with inherent limits
				\item \textbf{Risk Analysis:} Always uncertain, models approximate reality
				\item \textbf{Risk Evaluation:} Inherently subjective, needs calibration
				\item \textbf{BIA:} Dependencies compound recovery
				\item \textbf{Risk Treatment:} Optimize, cannot eliminate
				\item \textbf{Monitoring:} Continuous, but has limits
				\item \textbf{Reporting:} Actionable, audience-appropriate
			\end{itemize}
		\end{block}
	\end{frame}
	
	\begin{frame}{Final Wisdom}
		\begin{alertblock}{The Ultimate Truth}
			Risk management is not about eliminating risk.\\
			It is about understanding risk and making informed decisions.\\
			The goal is to optimize risk, not minimize it at all costs.\\
			\\
			\textbf{Risk cannot be eliminated, only managed.}\\
			\textbf{Complete risk identification is impossible.}\\
			\textbf{Risk analysis is always uncertain.}\\
			\textbf{But we must still decide.}
		\end{alertblock}
	\end{frame}
	
	\begin{frame}{Thank You}
		\begin{center}
			\Huge
			\textbf{Good Luck on Your CRISC Exam!}
			
			\vspace{1cm}
			
			\normalsize
			Remember: The ISACA mindset is about\\
			understanding why a question is asked,\\
			not just what is technically correct.
			
			\vspace{1cm}
			
			\textit{This completes 200 extremely difficult Domain B questions.}
		\end{center}
	\end{frame}
	
\end{document}